\title{Byte Latent Transformer: Patches Scale Better Than Tokens}

\begin{abstract}
We present the Byte Latent Transformer ({\textsc BLT}), a novel byte-level large language model architecture that, for the first time, achieves parity with token-based LLMs at scale while offering considerable gains in both inference robustness and efficiency.
In {\textsc BLT}, input bytes are grouped into dynamically determined patches, which function as the core computational entities. Patch boundaries are set according to the entropy of the upcoming byte, thereby directing computational and modeling resources to segments where the data is most complex. We provide the first computationally controlled scaling analysis of byte-level architectures, extending up to 8B parameters and 4T bytes of training data. Our findings verify that large-scale training on raw byte inputs—without a pre-defined vocabulary—is not only practical but advantageous. By adaptively utilizing longer patches when input sequences are easier to model, {\textsc BLT} achieves improved efficiency in both training and inference phases, while also exhibiting marked qualitative advances in reasoning and out-of-distribution generalization. Taken together, for equivalent inference expenditure, {\textsc BLT} surpasses tokenization-dependent approaches in scaling, benefiting from simultaneous increases in both model size and patch length.
\end{abstract}

\section{Introduction}
\label{section:intro}

We present the Byte Latent Transformer~(\textbf{BLT}), an architecture that operates directly on raw bytes without relying on tokenization, and demonstrates—for the first time—parity with tokenization-dependent models at large scale while offering notable gains in efficiency and robustness (\S\ref{sec:robustness}).
Current large language models (llms) are predominantly trained in an end-to-end fashion, with the notable exception of tokenization. Tokenization constitutes a heuristic, pre-processing phase that maps sequences of bytes to a predetermined set of tokens. This process imposes a particular compression schema on the data, resulting in challenges such as increased sensitivity to domain and modality~\citep{dagangetting}, pronounced vulnerability to input noise~(\S\ref{sec:robustness}), limited orthographic generalization~\citep{edman2024cute}, and persistent issues of multilingual inequality~\citep{liang2023xlm,petrov2024language,limisiewicz2024myte}.

Tokenization has traditionally played a crucial role, as training large language models on raw byte sequences incurs substantial computational expense, primarily due to increased sequence lengths~\citep{xue2022byt5}. Earlier studies have addressed this challenge by introducing more efficient self-attention mechanisms~\citep{el2020characterbert, clark2022canine} or by exploring architectures that eliminate attention entirely~\citep{wang2024mambabyte}~(\S\ref{section:related_work}). Nonetheless, these solutions are mainly effective for the training of \textit{smaller models}. When scaling up, the majority of computational effort in a Transformer comes from the extensive feed-forward network layers operating over each byte, rather than from the attention component itself.

We introduce a flexible, adaptive approach for partitioning bytes into \textit{patches}~(\S\ref{section:patching}), alongside a novel architecture that integrates both byte-level and patch-level signals. In contrast to tokenization methods, {\textsc BLT} does not rely on a static patch vocabulary. Instead, it utilizes compact, trainable encoder and decoder components to transform any subset of bytes into latent representations corresponding to patches. Our findings demonstrate that this approach yields a \textit{higher} efficiency in compute allocation compared to models based on tokenization.

Tokenization-based language models assign an equal computational budget to each token, prioritizing straightforward implementation over adaptive efficiency. However, the tokenization process employs compression heuristics that do not always align with the predictive difficulty of each token. Our architecture is grounded in the principle that computation should be flexibly distributed in accordance with demand. For instance, predicting word endings typically requires less computational effort—these are generally low-entropy, more predictable outcomes—whereas generating the initial word of a sentence poses a higher degree of uncertainty. This design philosophy is embodied in {\textsc BLT}'s structure~(\S\ref{section:architecture}), which incorporates three transformer modules: a pair of compact byte-level \textit{local models} and one expansive global \textit{latent transformer}~(\autoref{fig:arch_high_level}).

In order to determine optimal patch formation and the corresponding dynamic allocation of computation, {\textsc BLT} segments input data according to the entropy in the next-byte predictions. This approach yields contextual byte groupings characterized by comparably consistent information density.

We introduce the first systematic scaling investigation of byte-level models controlled for floating point operations (flops), covering model sizes up to 8B parameters and training on up to 4T bytes. Our findings demonstrate that end-to-end training directly from raw bytes at this scale is feasible, obviating the need for predefined-vocabulary tokenization. Importantly, {\textsc BLT} achieves comparable performance to Llama 3 when measured by training flops\footnote{The computational load of a model is quantified by tallying the Floating Point OPerations (flops) required.}, yet reduces inference time flops by as much as 50\%~(\S\ref{section:scaling}).

Moreover, leveraging raw byte inputs yields considerable advantages in capturing rare patterns within the data. Compared to tokenizer-driven models, {\textsc BLT} exhibits improved resilience to noisy text and superior competence in character-level phenomena—evident in tasks involving orthographic processing, phonological rules, and translation between low-resource languages~(\S\ref{sec:robustness}).

Additionally, {\textsc BLT} permits concurrent scaling of both model capacity and input patch size without exceeding a fixed inference flop threshold. Increasing patch lengths generally lowers overall computation per inference, freeing up capacity that can be used to expand the global latent transformer, since fewer passes are required. Our scaling trials constrained by inference flops~(\autoref{fig:fixed_inference_scaling}) reveal much more favorable scaling characteristics relative to traditional tokenization-based model designs.

To conclude, this work offers the following main contributions: 1) We present {\textsc BLT}, a byte latent llm framework designed to dynamically modulate computational resources in order to enhance flop efficiency; 2) We demonstrate that our approach reaches parity with Llama 3 in terms of flop-controlled training, up to the 8B parameter scale, while also enabling the possibility of sacrificing slight performance in evaluation metrics for up to 50\% improvements in flop efficiency; 3) {\textsc BLT} provides a novel scalability pathway for llms, allowing model capacity to grow while keeping inference cost constant; 4) We provide evidence of the superior resilience of {\textsc BLT} to noisy input, as well as its ability to capture sub-word structure within input data—features that are often overlooked by token-centric llms. The full training and inference codebase for {\textsc BLT} is publicly accessible at~\url{https://github.com/facebookresearch/blt}.

\section{Patching: From Individual Bytes to Groups of Bytes}
\label{section:patching}

Dividing byte sequences into distinct \textit{patches} enables {\textsc BLT} to flexibly assign computational resources according to contextual demands. As illustrated in Figure~\ref{fig:patching_types}, a variety of strategies exist for partitioning bytes into patches. More specifically, a patching function $f_p$ takes a byte sequence $\pmb{x}=\{x_i,|i=1,\ldots n\}$ of length $n$, and partitions it into $m < n$ patches $\pmb{p}=\{p_j|j=1,\ldots,m\}$ by assigning each $x_i$ to the set \{0,1\}, where a value of 1 marks the commencement of a new patch. In both token-based and patch-based approaches, the dominant computational expense results from the number of processing steps in the primary Transformer network. For {\textsc BLT}, this expense equates to the count of patches generated by a chosen patching function. Therefore, the mean size of a patch—referred to as \textit{patch size}—principally governs the data processing cost during both the training and evaluation phases for any particular patching scheme~(\S\ref{section:flops}).

In the sections that follow, we outline three methods for constructing patches: using a constant byte count per patch~(\S\ref{section:static-patch}), leveraging whitespace as patch boundaries~(\S\ref{section:white-patch}), and employing entropy signals from a compact byte-level language model to adaptively form patches~(\S\ref{section:dyn-patch}). We conclude by examining incremental patching and contrasting patching with the process of tokenization~(\S\ref{section:bpe-patch}).

\subsection{Strided Patching Every K Bytes}
\label{section:static-patch}

A commonly used approach for grouping bytes involves dividing them into fixed-size patches of length $k$, as utilized in MegaByte~\citep{yu2023megabyte}. This method's uniform stride simplifies both the training and inference pipelines and provides a direct way to adjust the mean patch size, thereby offering straightforward control over the computational cost. Despite these advantages, this strategy has notable limitations. Primarily, computational resources are not adaptively distributed according to need: for example, excessive computation may be spent predicting sequences with low informational content, like whitespace in code, while complex regions—such as mathematical expressions—may not receive adequate computational attention. Additionally, the use of static patching results in irregular and context-insensitive segmentation of similar byte sequences, which can manifest as the same word being split in differing ways depending on its position.

\subsection{Space Patching}
\label{section:white-patch}

According to \citet{slagle2024spacebyte}, an effective modification to standard strided patching involves generating new patches at every occurrence of a space-like byte\footnote{
    Space-like bytes refer to those that are not Latin characters, digits, or utf-8 continuation bytes. It is additionally required that each patch includes at least one non-space-like byte.}, leveraging these positions as natural segmentation points in many languages. In this Space patching approach, the latent transformer is invoked (consuming more computation) for each word boundary, thereby ensuring consistent patching of words throughout sequences and directing computational resources to challenging predictions, which often follow space-like characters. For instance, predicting the initial byte of the response to ``Who composed the Magic Flute?\uline{\hspace{.6em}}'' poses a greater challenge compared to subsequent bytes—once the first character is known (``M''), the possible continuations narrow significantly, making the remainder of ``Mozart'' relatively straightforward to predict. Nevertheless, Space patching exhibits limitations: it does not accommodate every language and domain equally well, and more crucially, lacks the capability to adjust patch sizes. In the following section, we describe a novel patching approach that exploits the observation that the first bytes of words tend to be the most challenging to anticipate, while also introducing a principled way to modulate patch size.

\subsection{Entropy Patching: Using Next-Byte Entropies from a Small Byte LM}
\label{section:dyn-patch}

Instead of utilizing traditional rule-based heuristics like whitespace to determine boundaries, our method employs a data-centric strategy to locate points of elevated uncertainty in next-byte predictions. We propose \textit{entropy patching}, a technique that leverages entropy measurements to inform the selection of patch boundaries.

To this end, we fit a compact byte-level auto-regressive language model on the {\textsc BLT} training set and estimate the entropy for predicting the next byte, considering the LM-generated distribution $p_{e}$ over the set of possible bytes $\mathcal{V}$:

\begin{equation}
H(x_i) = \sum_{v \in \mathcal{V}} p_{e}(x_i=v|\pmb{x}_{<i}) \log p_{e}(x_i=v|\pmb{x}_{<i})
\end{equation}

We explore two approaches for detecting patch boundaries using entropies $H(x_i)$. The first approach selects locations where the entropy exceeds a predefined global threshold, as shown in~\autoref{fig:patching}. The second method flags points where the entropy is notably higher compared to the preceding value. This latter technique can alternatively be understood as recognizing instances that disrupt the roughly monotonically decreasing entropy within a patch.

\begin{align*}
    \text{Global Constraint}\hspace{1cm}H(x_t)& > \theta_g\\
    \text{Approx. Monotonic Constraint}\hspace{1cm}H(x_t)& - H(x_{t-1}) > \theta_r
\end{align*}

Patch boundaries are determined through a simple preprocessing procedure performed at the data loading stage. This contrasts with the approach of~\citet{nawrot-etal-2023-efficient}, where a classifier is utilized to predict patch boundaries based on entropy. In our study~(\S\ref{section:experiments}), we evaluate and contrast these two strategies for separating bytes with low versus high entropy.

\subsection{The Byte-Pair Encoding (BPE) Tokenizer and Incremental Patching}
\label{section:incremental-patching}
\label{section:bpe-patch}

A large number of current llms, such as our baseline Llama 3, rely on subword tokenization methods like bpe~\citep{gage1994new, sennrich-etal-2016-neural}. Throughout this work, we use the term ``tokens'' to indicate byte sequences selected from a \textit{finite} vocabulary that is fixed before the training phase, distinguishing them from ``patches,'' which are composed of adaptively grouped sequences that are not constrained by a predetermined vocabulary. One fundamental distinction is that models using tokens are unable to directly access the original byte-level characteristics, unlike when working with patches.

An important advancement offered by {\textsc BLT} compared to models that rely on tokenization lies in its reconfiguration of the balance between vocabulary size and computational requirements. In conventional large language models, expanding the vocabulary typically results in longer average token lengths, leading to a reduction in the number of processing steps but at the expense of an enlarged output dimension in the model's final projection layer. This balance restricts tokenization-driven models from effectively optimizing both token size and inference efficiency. As an illustration, Llama 3 achieves an increase in the mean token length from 3.7 to 4.4 bytes; however, this comes at the cost of enlarging its embedding table by a factor of four when compared to Llama 2.

During generation, {\textsc BLT} must determine at each point in the byte sequence whether it is positioned at a patch boundary, as this dictates whether additional computation involving the Latent Transformer is required. This assessment must be made in isolation, without reference to the remainder of the sequence that has yet to be produced. Consequently, the process of patching must be agnostic to any future bytes, and cannot utilize information from later in the sequence to inform its segmentation choices. More formally, a patching scheme $f_p$ is incrementally patching if it adheres to the following property:
$$f_p(\pmb{x}_{<i}) = f_p(\pmb{x})_{<i}$$
bpe does not constitute an incremental patching scheme, because the same initial substring can be segmented differently depending on what follows, and thus violates this criterion\footnote{The introduction of a dedicated delimiter token to indicate patch boundaries can make bpe incremental, but this comes at the cost of elongating the byte sequence.}.

\section{{\textsc BLT} Architecture}
\label{section:architecture}
{\textsc BLT} comprises a large, globally autoregressive language model that processes patch-level representations, coupled with two lightweight local models: one responsible for encoding contiguous bytes into patches and the other for reconstructing bytes from patch representations~(\autoref{fig:arch_high_level}).

\subsection{Latent Global Transformer Model}

\textit{The Latent Global Transformer} refers to an autoregressive transformer model, denoted as $\mathcal{G}$ and composed of $l_{\mathcal{G}}$ layers. This model processes an ordered sequence of latent input patch representations $p_j$ and produces a corresponding sequence of output patch representations $o_j$. In our notation, patch indices are consistently represented by the subscript $j$, whereas byte indices are indicated by $i$. The global model applies a block-causal attention mask~\citep{dubey2024llama}, which enforces that each patch may only attend to itself and all preceding patches within the given document. Notably, this component dominates computational cost during both pre-training and inference. As a consequence, by deciding when to apply this model, one can modulate the computational resources spent on particular segments of the input or output in accordance with their complexity.

\subsection{Local Encoder}
\label{section:local}

\textit{The Local Encoder Model}, indicated by $\mathcal{E}$, is designed as a compact, transformer-based encoder, comprised of $l_{\mathcal{E}} << l_{\mathcal{G}}$ layers. Its central purpose is to rapidly convert an input stream of bytes $b_i$ into a set of rich patch embeddings, $p_j$. Distinct from conventional transformer architectures, each transformer layer is immediately followed by a cross-attention layer, which serves to aggregate byte-level features into patch representations~(\autoref{fig:crossattn}).

To begin, the raw byte sequence $b_i$ is projected into embedding vectors $x_i$ via a learned embedding matrix of size $\mathbb{R}^{256 \times h_{\mathcal{E}}}$. Optionally, these embeddings may incorporate supplementary information by means of hash-embeddings~(\S\ref{sec:hash-emb}). This representation is subsequently refined through an interleaved stack of transformer and cross-attention layers~(\S\ref{sec:enc-cross-attn}), yielding patch-level representations $p_i$, which are then passed to the higher-level global transformer $\mathcal{G}$ for further computation.

The attention mechanism employed in these transformer layers is a \textit{local block causal} scheme: each position can attend to a bounded window of $w_{\mathcal{E}}$ previous bytes, permitting attention across dynamic patch divisions but disallowing it past document boundaries. The following sections offer further insight into the embedding process and the design of the cross-attention module.

\subsubsection{Encoder Hash n-gram Embeddings}
\label{sec:ngram-emb}
\label{sec:hash-emb}
An essential factor in developing rich and stable representations at each timestep $i$ involves encoding details about prior bytes. In {\textsc BLT}, this is accomplished by representing both the current byte $b_i$ on its own as well as within the context of its corresponding byte n-gram. Specifically, at every position $i$, byte-grams are first generated

\begin{align}
    g_{i,n}=\{b_{i-n + 1},\ldots, b_i\}
\end{align}

for each byte position $i$ and $n$ from three to eight.\footnote{
We omit byte-grams of size $n$ or more when $i<n$. }

Next, we present hash $n$-gram embeddings, where each byte $n$-gram is assigned to an index in a fixed-size embedding table $E_{n}^{hash}$ using a hash function, applied independently for each $n \in \{3,4,5,6,7,8\}$~\citep{bai2010learning}. The embedding derived from this mapping is subsequently summed with the corresponding byte embedding, after which it undergoes normalization and is provided as input to the local encoder model. We compute the augmented embedding

\begin{align}
e_i &= x_i + \sum_{n = 3,...,8}E_{n}^{hash}(\text{Hash}(g_{i,n})) \\
\text{where, Hash}(g_{i,n}) &= \text{RollPolyHash}(g_{i,n}) \% |E_{n}^{hash}|
\end{align}

We scale $e_i$ by dividing by the total number of distinct $n$-gram sizes, incremented by one, and apply the RollPolyHash function as described in Appendix~\ref{appendix:rollpolyhash}. Section~\ref{section:ablations} presents a study of how varying both the size of $n$ and the embedding table impacts the scaling law dynamics when controlling for computational budget, focusing specifically on $n$-gram hash embeddings. Alongside these hash-based embeddings, we also conducted experiments utilizing frequency-based $n$-gram embeddings; further information regarding these investigations can be found in Appendix~\ref{appendix:ngrams}.

\subsubsection{Encoder Multi-Headed Cross-Attention}
\label{sec:enc-cross-attn}
Our approach adheres closely to the input cross-attention mechanism established in the Perceiver framework~\citep{jaegle2021perceiver}, but with an important distinction: here, the latent representations are aligned with patch-specific representations rather than a predetermined latent set~(\autoref{fig:crossattn}), and attention is limited solely to the byte elements comprising each patch. The cross-attention module utilizes a query vector per patch $p_j$, initially generated by applying a pooling operation over the byte-level representations associated with patch $p_j$, followed by a linear transformation via $\mathcal{E}_{C} \in \mathbb{R}^{h_{\mathcal{E}} \times (h_{\mathcal{E}} \times U_{\mathcal{E}})}$, with $U_{\mathcal{E}}$ indicating the number of attention heads used in the encoder cross-attention. Explicitly, if $f_{\text{bytes}}(p_j)$ maps to the byte sequence within patch $p_j$, computation proceeds as follows:

\begin{align}
P_{0,j} &= \mathcal{E}_{C}(f_\text{bytes}(p_j)), \quad f~ \text{denotes a pooling operation} \\
P_l &= P_{l-1} + W_o\left(\text{softmax}\left(\frac{QK^T}{\sqrt{d_k}}\right)V\right) \\
\text{with } Q_j &= W_q(P_{l-1,j}), \;\; K_i = W_k(h_{l-1,i}), \;\; V_i = W_v(h_{l-1,i}) \\
h_l &= \text{Encoder-Transformer-Layer}_l(h_{l-1})
\end{align}

Here, $P \in \mathbb{R}^{n_p \times h_{\mathcal{G}}}$ denotes the $n_p$ patch representations fed into the global model. These representations are created by aggregating the byte embeddings $e_i$ corresponding to each individual patch $p_j$. The matrices $W_q$, $W_k$, $W_v$, and $W_o$ serve as the projection weights for queries, keys, values, and the output, respectively, where keys and values are obtained by projecting the byte-level representations $h_i$ generated from the preceding layer (or $e_i$ for the initial layer). A specialized masking approach for patching is applied such that each query $Q_j$ is permitted to attend only to those keys and values corresponding to the bytes within its respective patch $j$. Given that multi-head attention is applied to $Q, K,$ and $V$ and that patch representations generally occupy a higher dimensionality ($h_{\mathcal{G}}$) than $h_{\mathcal{E}}$, $P_l$ is maintained as multiple heads each with dimension $h_{\mathcal{E}}$ during cross-attention; subsequently, these are concatenated into a vector of dimensionality $h_{\mathcal{G}}$. The module also includes a pre-LayerNorm operation on queries, keys, and values, and intentionally omits positional embeddings in the cross-attention mechanism. Lastly, a residual connection wraps the cross-attention block to facilitate optimization.

\subsection{Local Decoder} 

Analogous to the architecture of the local encoder, the local decoder $\mathcal{D}$ utilizes a compact transformer-based design characterized by $l_{\mathcal{D}} << l_{\mathcal{G}}$ layers. This component reconstructs raw bytes $y_i$ from the sequence of global patch embeddings $o_j$. During decoding, it generates each raw byte conditioned on the bytes decoded up to that point, leveraging hidden states from the local encoder corresponding to the byte sequence. The local decoder stacks $l_{\mathcal{D}}$ layers that alternate between cross-attention mechanisms and transformer blocks. Specifically, cross-attention precedes the transformer module in each layer to map patch-based representations into byte-level representations, after which the transformer block processes these byte representations sequentially.

\subsubsection{Decoder Multi-headed Cross-Attention} 

Within the decoder's cross-attention module, the typical configuration of queries and key/value pairs is reversed; that is, byte-representations act as queries while patch representations provide the keys and values. To initiate the cross-attention computation, the byte-representations are seeded with the embeddings produced by the last layer of the encoder, denoted $h_{l_{\mathcal{E}}}$. For each subsequent decoder layer $l$, the byte-representations $d_{l,i}$ are determined by the following procedure:

\begin{align}
D_0 &= h_{l_{\mathcal{E}}} \\
B_l &= D_{l-1} + W_o\left(\text{softmax}\left(\frac{QK^T}{\sqrt{d_k}}\right)V\right), \\
  &\text{where } Q_i = W_q(d_{l-1,i}), \quad K_i = W_k(\mathcal{D}_C(o_j)), \quad V_i = W_v(\mathcal{D}_C(o_j)) \\
  D_l &= \text{Decoder-Transformer-layer}_l(B_l)
\end{align}

In this context, $W_k$ and $W_v$ again serve as the projection matrices for keys and values, applying linear transformations combined with the split operation $\mathcal{D}_C$ to the terminal patch representations $o_j$ derived from the global model. The matrix $W_q$ acts as a query projection on the byte-level representations $d_{l-1}$ output by the preceding decoder transformer layer (or $h_{l_{\mathcal{E}}}$ at the initial layer). Meanwhile, $W_o$ projects outputs, ensuring that $B$ occupies $\mathbb{R}^{h_{\mathcal{D}} \times n_b}$, where $n_b$ designates the total number of output bytes. The updated decoder hidden states $D_l$ are obtained by passing the output of the cross-attention module, $B$, through an additional decoder transformer layer. Consistent with the approach used for local encoder cross-attention, our setup incorporates multi-head attention, employs pre-LayerNorm structures, omits positional embeddings, and introduces a residual pathway surrounding the cross-attention layer.

\section{Experimental Setup}
\label{section:experiments}
We establish a suite of controlled experiments to rigorously evaluate {\textsc BLT} in comparison with models utilizing tokenization. Special care is taken to ensure that {\textsc BLT} does not benefit from potentially increased sequence length, thereby enabling fair comparisons.

\subsection{Pre-training Datasets} Across all model sizes discussed in this work, pre-training utilizes two primary datasets: 1) the Llama 2 dataset~\citep{touvron2023llama}, a corpus totaling 2 trillion tokens sourced from diverse publicly available repositories, then subjected to rigorous cleaning and filtering procedures aimed at enhancing overall data quality; and 2) {\textsc BLT}-1T, a novel collection consisting of 1 trillion tokens aggregated from multiple public origins, which also incorporates a segment of the pre-training material published by Datacomp-LM~\citep{li2024datacomp}. The Llama 2 dataset primarily serves scaling law analyses—specifically, examining the optimal token count as proposed by~\cite{dubey2024llama}—to inform architectural decisions in the construction of {\textsc BLT}. In contrast, {\textsc BLT}-1T is utilized for full pre-training cycles designed to facilitate head-to-head evaluation against Llama 3 on downstream benchmarks. Notably, both datasets are curated without inclusion of any data from Meta’s products or services. In the context of baseline comparisons using tokenizer-based models, we rely on the Llama 3 tokenizer (featuring a 128K token vocabulary), which in our assessments delivers better performance than the tokenizer associated with Llama 2.

\subsection{Entropy Model} In our study, the entropy model operates at the byte level and is trained on the identical dataset as the comprehensive {\textsc BLT} model. Unless specified otherwise, our default architecture is a transformer consisting of 100 million parameters, arranged across 14 layers, with each hidden layer having a dimensionality of 512. The attention mechanism utilizes a sliding window over 512 bytes. All other hyperparameter settings are consistent with those of our local and global transformer variants. We explored a variety of model configurations, such as changing model capacity, receptive field sizes, and architectural choices, which are elaborated in \autoref{section:ablations}. Notably, when the model’s receptive field is sufficiently narrow, its trained entropy representations can be stored compactly as a lookup table.

\subsection{Entropy Threshold and Equalizing Context Length} When implementing entropy-based patching, we determine an entropy threshold that results in a target average \textit{patch size} across the pretraining dataset. In the case of {\textsc BLT}, contrary to traditional tokenization, the \textit{patch size} is freely configurable, which meaningfully affects the effective context length processed by the model. To standardize the average context and prevent any biases arising from using larger patch sizes, we set the expected number of bytes per batch to be consistent. Consequently, for configurations with larger patch sizes, we correspondingly shorten sequence lengths. Specifically, for Llama 2 data, the context window is set to 8k bytes, whereas for the {\textsc BLT}-1T dataset, the average context window is expanded to 16k bytes. Throughout, the average batch size is kept fixed at 16M bytes.

Although the mean batch size does not change, dynamic patching methods can produce varying bytes-to-patches ratios in the batches. To optimize efficiency, our {\textsc BLT} training implementation organizes batches directly by patches, eliminating the need for padding within the more computationally demanding latent transformer. As a result, all batches are standardized to contain an equal number of patches. For training purposes, byte sequences are padded or, if necessary, truncated to 12k bytes for Llama 2 and 24k bytes for the {\textsc BLT}-1T datasets. This practice helps prevent excessive memory consumption due to sequences that may contain abnormally large patches.

\subsection{Entropy Model Context}
\label{section:entropy-context}
Our empirical observations indicate that entropy patching tends to create increasingly larger patches within structured datasets, such as in multiple choice problems (refer to the example patching in an MMLU instance in \autoref{fig:mmlu_patching}), which frequently exhibit high repetitiveness. These discrepancies can be attributed to reduced entropy arising from the recurring content embedded in the entropy model's context. Therefore, for the full-scale application of {\textsc BLT}-Entropy with a patch size of 4.5, we mitigate these effects by refreshing the entropy context using newline characters and adopting an approximate monotonicity constraint—this alternative is less susceptible to "entropy drift" that may result from fluctuating context lengths. This adjustment solely influences entropy computations; the approach for determining the entropy threshold remains unchanged.

\subsection{FLOPs Estimation} 
\label{section:flops}

Our methodology for calculating transformer FLOPs closely mirrors that outlined in Chinchilla~\citep{hoffmann2022training}, encompassing FLOPs from the feed-forward network, qkvo projections within the self-attention mechanism, as well as attention computation and output projection stages. A key distinction in our approach is the assumption that the input embedding layer utilizes an efficient lookup mechanism rather than relying on a dense matrix multiplication, rendering it a FLOP-free operation. Consistent with established literature, we approximate the computational cost of the backward pass as being twice that of the forward pass.

To determine the number of flops \textit{per byte} for {\textsc BLT} architectures, we accumulate the flops incurred by the local encoder transformer, the global latent transformer, and the local decoder transformer, in addition to the contributions from the cross attention modules present in both the encoder and the decoder:

\begin{align}
    \text{FL}_{\text{{\textsc BLT}}} &= \text{Transf. FL}(h_{\mathcal{G}}, l_{\mathcal{G}}, m=n_{ctx}/n_p, V=0) / n_p\\
   &+\text{Transf. FL}(h_{\mathcal{E}}, l_{\mathcal{E}}, m=w_{\mathcal{E}}, V=0) \\
   &+ \text{Transf. FL}(h_{\mathcal{D}}, l_{\mathcal{D}}, m=w_{\mathcal{D}}, V=256) \\ 
    &+ \text{Cross Attn. FL}(h_{\mathcal{E}}, l_{\mathcal{E}}, m=n_p, r=n_p/k) \times k/n_p \\ 
   &+ \text{Cross Attn. FL}(h_{\mathcal{D}}, l_{\mathcal{D}}, m=k, r=k/n_p) 
\end{align}

Here, $n_{ctx}$ denotes the sequence length measured in bytes, and $n_p$ specifies the patch size. The parameter $r$ refers to the proportion of queries relative to key/values. The variable $k$ indicates the ratio between the patch dimension and the byte dimension; that is, it represents how many local model splits are joined together to yield the overall model representation ($k=2$ as shown in \autoref{fig:crossattn}). $V$ indicates the size of the output projection vocabulary, utilized solely within the local decoder. The length of the context for attention, denoted as $m$, varies depending on whether a module processes the byte or patch sequence. We adjust the calculation of attention flops for each module in accordance with this distinction. Full mathematical formulations for the computation of Transformer-FLOPs and Cross-Attention FLOPs can be found in Appendix~\ref{section:flops-appendix}.

\subsection{Bits-Per-Byte Estimation} Perplexity is meaningful only when the tokenizer remains consistent, since it quantifies the model’s uncertainty at the token level. To fairly evaluate and compare models that operate at both byte and token levels, prior research~\citep{xue2022byt5, yu2023megabyte, wang2024mambabyte} advocates for reporting Bits-Per-Byte (BPB). BPB serves as a variant of perplexity that does not depend on the choice of tokenizer. In particular:

\begin{align}
    \text{BPB}(x)&= \frac{\mathcal{L}_{CE}(\pmb{x})}{\ln(2)\cdot n_{\text{bytes}}}
\end{align}

where the uncertainty over the data $\pmb{x}$ as measured by the sum of the cross-entropy loss is normalized by the total number of bytes in $\pmb{x}$ and a constant. 

\subsection{Transformer Architecture Hyperparameters} The transformer blocks utilized in {\textsc BLT}—including both the local and global components—are modeled primarily after the Llama~3~\citep{dubey2024llama} configuration. Specifically, the feed-forward networks employ the SwiGLU activation function~\citep{swiglu}, while rotary positional embeddings (RoPE)~\citep{RoPe} with $\theta=500000$~\citep{xiong2024effective} are incorporated exclusively within the self-attention modules. RMSNorm \citep{RMSNorm} is adopted for all layer normalization stages. Self-attention layers that utilize fixed-standard masking schemes (such as \textit{block causal} or \textit{fixed-window block causal}) are implemented using Flash attention~\citep{dao2022flashattention}, with a window size of 512 applied for masks of fixed width. For the cross-attention layers, which require dynamic, patch-dependent masking, Flex Attention\footnote{\url{https://pytorch.org/blog/flexattention}} is employed; this provides optimized fused kernels, thereby accelerating the training process considerably.

\subsection{{\textsc BLT}-Specific Hyperparameters} To evaluate the capabilities of {\textsc BLT} models, we carry out experiments focusing on both scaling behavior and performance on downstream tasks, assessing models with sizes of 400M, 1B, 2B, 4B, and 8B parameters. Detailed configuration hyperparameters for these architectures are available in Appendix~Table~\ref{tab:arch}.

Initialization of queries for the initial cross-attention layer in the local encoder is performed using max-pooling. Consistently across all {\textsc BLT} variants, we employ $500,000$ hashes using a single hash function and n-gram sizes set between 3 and 8. A fixed learning rate of $4e-4$ is adopted throughout. Our decision to use a matched learning rate for both token-based and {\textsc BLT} models results from hyperparameter sweeps between $1e-3$ and $1e-4$ at the 400M and 1B model scales, which demonstrated that this shared value yields optimal outcomes. Training batch sizes for scaling experiments on Llama-2 data adhere to guidelines suggested by~\cite{dubey2024llama}, matched for data volume in bytes where appropriate. Optimization is managed using AdamW~\citep{adamw} with $\beta_1$ set at 0.9, $\beta_2$ at 0.95, and $\epsilon=10^{-8}$. A linear warm-up of 2000 steps is employed, after which the learning rate follows a cosine decay to zero. Additionally, we apply weight decay with a factor of 0.1, and clip the global gradient norm at 1.0.

\section{Scaling Trends}
\label{section:scaling}

We offer a comprehensive overview of scaling behaviors exhibited by byte-level models, with the intention of guiding future expansion of {\textsc BLT} models. Our scaling investigation seeks to overcome gaps present in earlier work on byte-level modeling through the following contributions: (a) we analyze scaling patterns within a compute-optimal training framework, (b) we construct and evaluate 8B parameter models using substantial training corpora—reaching up to 1T tokens (4T bytes)—and assess their efficacy on various downstream applications, and (c) we assess how scaling manifests when inference costs are held constant. Subsequent sections will further explore particular benefits that arise from processing byte-sequences directly.

\subsection{Parameter Matched Compute Optimal Scaling Trends}
\label{section:parameter-matched-scaling}
For this study, we utilize the Llama 2 dataset to train a range of \textit{compute-optimal} bpe and {\textsc BLT} models spanning four scales, with parameter counts from 1B to 8B. We examine the relationship between the training flops and language modeling quality, evaluating these models on a representative portion of the training data mix. The bpe models follow the parameter-to-token ratio identified as optimal by Llama~3~\citep{dubey2024llama}. Such a \textit{compute-optimal} configuration is intended, from a theoretical perspective, to maximize training set performance within a fixed computational budget~\citep{hoffmann2022training}, serving as a robust benchmark for comparison. Each bpe model is paired with a {\textsc BLT} model, which is trained on identical data and employs a Latent Transformer of the same architecture and scale.

As shown in~\autoref{fig:scaling-compute-opt} (right), {\textsc BLT} models consistently equal or surpass the performance of their bpe-based peers, a pattern that persists across increases in both model size and computational resources. To our knowledge, {\textsc BLT} represents the first byte-level Transformer design to demonstrate comparable scaling behaviors with BPE-based systems in regimes optimized for compute. This observation supports our hypothesis that the ideal parameters-to-compute ratio for bpe models is similarly applicable to {\textsc BLT}, or at least remains closely aligned.

Both enhancements to model architectures and the integration of dynamic patching mechanisms are essential to align with bpe scaling trends. As depicted in ~\autoref{fig:scaling-compute-opt} (left), we assess models that utilize space-patching, directly comparing them to Llama 3. Here, SpaceByte \citep{slagle2024spacebyte} is approximated using {\textsc BLT} configured for space-patching but omitting n-gram embeddings and cross-attention components. While SpaceByte demonstrates gains relative to Megabyte, its performance remains substantially lower than that of Llama 3. The right panel of \autoref{fig:scaling-compute-opt} showcases the respective contributions from both improvements to architecture and the application of dynamic patching. At large scales, {\textsc BLT} models achieve results commensurate with leading tokenizer-based models, exemplified by Llama 3.

Furthermore, we note the impact of tokenizer selection on tokenizer-based model performance: when training on the same dataset, models employing the Llama-3 tokenizer surpass those utilizing the Llama-2 tokenizer.

In summary, when employing much larger patch sizes, our {\textsc BLT} framework demonstrates performance characteristics that lie between those of Llama 2 and Llama 3. While the average token sizes for Llama 2 and Llama 3’s bpe tokenizers are 3.7 and 4.4 bytes, respectively, {\textsc BLT} is able to mirror these scaling behaviors with substantially larger average patch sizes—specifically, 6 and even 8 bytes. Because inference flops are inversely related to patch size, selecting a patch size of 8 bytes yields close to a 50\% reduction in inference flops. Additionally, increasing patch size appears to enhance model performance as both model capacity and dataset size expand. Notably, while {\textsc BLT} with a patch size of 8 begins at a much lower performance than bpe Llama 2 at the 1B scale, it ultimately surpasses bpe performance at the 7B scale. This observation implies that larger patch sizes may further improve as model and data scales continue to grow, and even larger patch sizes might become practical with advancements in model scale and training computation.

\subsection{Beyond Compute Optimal Task Evaluations}
\label{section:evals}
To further examine how scaling impacts performance, we train an 8B {\textsc BLT} model using more compute than the compute optimal ratio on the {\textsc BLT}-1T dataset, a larger and higher-quality corpus. Model capabilities are then evaluated using a variety of established classification and generation benchmarks that test reasoning, general knowledge, and coding skills.

\paragraph{Classification tasks} considered in our evaluation include ARC-Easy (0-shot)~\citep{Eval_ARC}, Arc-Challenge (0-shot)~\citep{Eval_ARC}, HellaSwag (0-shot)~\citep{Eval_hellaswag}, PIQA (0-shot)~\citep{Eval_piqa}, and MMLU (5-shot)~\citep{hendrycks2020measuring}. For these tasks, we apply prompt-scoring, which involves calculating the probability assigned to each choice character, and aggregate results as mean accuracy.

\paragraph{Coding related generation tasks:} To measure the model's Python code generation proficiency, we present pass@1 scores on MBPP (3-shot)~\citep{Eval_MBPP} and HumanEval (0-shot)~\citep{Eval_HumanEval}.

In \autoref{tab:evals}, we present a comparison among three models trained on the {\textsc BLT}-1T dataset: a bpe Llama 3 tokenizer-based model,\footnote{We select the Llama 3 tokenizer with a vocabulary size of 128k, given its superior performance over the Llama 2 tokenizer's 32k vocabulary.} as well as two versions of the {\textsc BLT} model. These include a space-patching variant ({\textsc BLT}-Space) and an entropy-based patching variant ({\textsc BLT}-Entropy). The latter is trained with an approximate monotonicity constraint, and its context is reset with new lines as detailed in~\autoref{section:entropy-context}. Each model is allocated the same flop budget for training. Notably, for {\textsc BLT}-Entropy, we lower the inference entropy threshold from 0.6 to 0.1, which enhances performance on evaluated tasks but necessitates a greater number of inference steps.

The {\textsc BLT}-Entropy model achieves superior performance compared to the Llama 3 model in 4 out of 7 evaluated tasks, despite both being trained with an equivalent number of bytes. This enhanced performance is likely attributable to two primary factors: (1) more efficient utilization of training compute enabled by dynamic patching, and (2) the explicit modeling of data at the byte level rather than at the token level.

Conversely, {\textsc BLT}-Space performs worse than the Llama 3 tokenizer across nearly all tasks, except for one; however, its adoption of a larger average patch size of 6 bytes allows it to substantially lower inference flops. By contrast, both the bpe and entropy-patching models display similar average patch sizes, each around 4.5 bytes within the mixed training dataset. Given a fixed training budget, the model with the larger patch size is able to process 30\% more data compared to the other two approaches. This broader data coverage may shift {\textsc BLT} even further from the point of compute optimality.

\subsection{Patches Scale Better Than Tokens} In {\textsc BLT} models, it is possible to scale up both the model capacity and the patch size in tandem, all while preserving a fixed budget for training and inference computations, as well as holding the training dataset size constant. Unlike fixed-vocabulary token-based approaches, patch-based architectures uniquely allow arbitrary increases in patch size, overcoming the efficiency limitations inherent in tokenization schemes, as detailed in Section~\ref{section:bpe-patch}. Using longer patches reduces computational requirements, which can then be redirected to expand the global latent transformer, since it needs to be invoked less frequently.

In order to evaluate whether larger models executing fewer steps on bigger patches outperform smaller counterparts that take more steps, we carry out a fixed inference scaling analysis. Using initial model sizes of 400m and 3.6B parameters paired with the Llama 2 tokenizer, we identify flop-matched models utilizing the Llama 3 tokenizer and {\textsc BLT}-Entropy architectures with mean patch sizes set at 6 and 8 bytes, measured on the training datamix (refer to~\autoref{tab:fixed-inf-params} for specific model configurations). Notably, for models configured with an 8-byte patch size, three encoder layers are adopted in place of a single layer. All models are trained across varying training flop allocations.

As illustrated in \autoref{fig:fixed_inference_scaling}, {\textsc BLT} models demonstrate superior scaling behaviors in comparison to tokenization-based architectures across both inference flop categories. Initially, BPE models exhibit better performance when training resources are limited, but {\textsc BLT} rapidly overtakes them, surpassing the compute-optimal threshold. In real-world scenarios, investing additional resources in one-time pretraining can be advantageous, as it enables the development of higher-quality models while maintaining a fixed inference computational cost. A notable case is the 8B parameter model class, including Llama 3.1, which was pretrained on a dataset two orders of magnitude larger than what would be considered compute-optimal for models of its size.

The threshold at which {\textsc BLT} surpasses token-based models has moved marginally closer to the compute-optimal point with the transition to models in a higher flop category—shifting from requiring 3 times the compute-optimal budget to just 2.5 times. Additionally, the model with patch size 8 demonstrates a more pronounced scaling advantage in the higher flop class, allowing it to outperform the other models at an earlier stage. As detailed in Section~\ref{section:parameter-matched-scaling}, increasing patch sizes tends to bring performance more in line with BPE models as model scale increases. We partially attribute this phenomenon to the smaller proportion of total flops consumed by the byte-level Encoder and Decoder, which scale up less rapidly than the Latent Transformer. Notably, scaling the total parameter count twentyfold, from 400M to 8B, results in only about a twofold increase in {\textsc BLT}'s local parameters. This detail is significant, because larger patch sizes impact the flops solely in the patch Latent Transformer, leaving the byte-level components unchanged. This mechanism explains why, with larger models, the {\textsc BLT}-Entropy configuration with ps=8 increased only modestly from 1.6x to 1.7x relative to the Llama 2 model size.

In conclusion, our investigation into patch-length scaling reveals that the {\textsc BLT} patch-based architecture exhibits improved scaling behavior when both patch size and model capacity are increased together. Notably, this positive trend continues and may even be enhanced as the model becomes larger.

\section{Byte Modeling Improves Robustness} Additionally, we evaluate the robustness of {\textsc BLT} relative to token-based models which do not incorporate explicit byte-level data, and outline a method for transforming pretrained token-based models into a byte-oriented format.
\label{sec:robustness}

\subsection{Character-Level Tasks}

One initial impetus for developing byte-level models was their inherent resilience to noise at the byte level within inputs, as well as their ability to recognize the internal structure of tokens—a feature often lacking in contemporary tokenizer-based models. To assess these particular capabilities, we conduct further experiments using benchmarks specifically designed to test both input noise robustness and character-level understanding, encompassing English, various other languages, numerals, and phonemes. The outcomes of these assessments are summarized in Table~\ref{tab:char_tasks}.

\paragraph{Noisy Data} To assess the robustness of {\textsc BLT} relative to tokenizer-based architectures, we generate altered forms of the standard classification datasets discussed in Section~\ref{section:evals}. We utilize five unique character-level perturbation techniques to introduce textual noise: (a) \textit{AntSpeak}, which converts all text to uppercase and inserts spaces between individual characters; (b) \textit{Drop}, which randomly deletes 10\% of characters from each instance; (c) \textit{RandomCase}, which randomly sets half the characters to uppercase and the other half to lowercase; (d) \textit{Repeat}, which repeats up to 20\% of characters (with a cap of four repetitions per character); and (e) \textit{UpperCase}, where all text is capitalized. For our evaluations, each perturbation method is applied independently to the prompt, the completion, or both—defining distinct experimental conditions. We aggregate performance across these configurations. Table \ref{tab:char_tasks} presents results for perturbed versions of HellaSwag~\citep{Eval_hellaswag}, demonstrating that {\textsc BLT} consistently surpasses tokenizer-based baselines in robustness, exhibiting an average performance improvement of 8 points compared to a similarly trained model, and even outperforming Llama 3.1 despite the latter's considerably larger training set.

\paragraph{Phonology - Grapheme-to-Phoneme (G2P)} We evaluate {\textsc BLT}'s proficiency in converting sequences of graphemes (individual characters denoting a word) into corresponding phonemic transcriptions (representations of pronunciation). Table \ref{tab:char_tasks} shows the 5-shot G2P task results on Phonology Bench~\citep{suvarna2024phonologybench}, revealing that {\textsc BLT} achieves superior performance compared to the baseline Llama 3 1T tokenizer-based model for this task.

\paragraph{CUTE}
In order to evaluate character-level comprehension, we test {\textsc BLT} on the CUTE benchmark~\citep{edman2024cute}, which is composed of multiple tasks generally grouped into three main areas: understanding composition, grasping orthographic similarity, and manipulating sequences. The CUTE benchmark presents a substantial obstacle for the majority of tokenizer-based systems; these models typically show awareness of their tokens’ internal spelling but find it difficult to leverage this for effective text manipulation. As seen in Table~\ref{tab:char_tasks}, {\textsc BLT}-Entropy achieves over 25 points higher performance than the BPE-based Llama 3 models on this suite. Of note, our approach excels particularly in tasks that demand character manipulation, achieving 99.9\% accuracy across both spelling-related tasks. These substantial gains, despite {\textsc BLT} being trained with merely 1/16th of the data of Llama 3.1, suggest that BPE tokenization frameworks struggle to capture character-level structure effectively. Figure \ref{fig:cute_examples} provides representative examples in which the Llama 3 tokenizer model fails, while our {\textsc BLT} system succeeds. The only areas where BPE demonstrates superior performance are word insertion and deletion tasks; although such manipulations may not be as natural for models that operate at the byte level, the difference is modest, and it may be more feasible to compose words from character-level information than to deconstruct tokens. All evaluations employ the same experimental setup and make use of the original Huggingface prompts. Additional prompt refinement could potentially enhance the performance of BPE-based models.

\paragraph{Low Resource Machine Translation} We assess the performance of {\textsc BLT} in translation tasks involving both directions for six major language families and twenty-one low-resource languages written in a variety of scripts, utilizing the FLORES-101 benchmark~\citep{goyal2022flores}. SentencePiece BLEU scores for these evaluations are presented in Table~\ref{tab:flores}. 
The findings indicate that {\textsc BLT} surpasses a counterpart model based on the Llama 3 tokenizer, providing an overall BLEU improvement of 2 points when translating into English, and a 0.5-point gain when translating out of English. While results for widely-used language pairs show that {\textsc BLT} matches or slightly exceeds Llama 3’s performance, the model demonstrates markedly superior results across numerous low-resource language pairs, particularly within less-resourced language families. This highlights the capacity of byte-level modeling in {\textsc BLT} to robustly handle infrequent and diverse byte sequences.

\subsection{Training {\textsc BLT} from Llama 3}
\label{sec:distillation}
We investigate a procedure in which {\textsc BLT} models capitalize on prior knowledge from pre-trained tokenizer-based architectures to facilitate more efficient and rapid convergence. This is done by transferring the global transformer weights from a pre-trained Llama 3.1 model to initialize the corresponding parameters in {\textsc BLT}. During training, the global transformer’s weights are updated with a learning rate set to one-tenth of that used for the local encoder and local decoder components, following the recommended training steps for Llama 3. Table~\ref{tab:distill} illustrates a comparative analysis against a standard {\textsc BLT} model. The results clearly show that initializing {\textsc BLT} with weights from Llama 3.1 leads to markedly better performance than both the standalone Llama 3 and baseline {\textsc BLT} models trained under equivalent compute budgets. Furthermore, when juxtaposed with our {\textsc BLT}-Entropy model from Table~\ref{tab:evals}—which utilizes a much larger dataset (1T tokens)—the Llama 3.1-initialized {\textsc BLT} continues to demonstrate improved results on the MMLU benchmark. This outcome suggests the method holds substantial promise for minimizing the total training FLOPs required.

This approach can be interpreted as adapting models that rely on tokenizers into ones that function without them, thereby transforming a pre-trained LLaMA 3.1 into a {\textsc BLT} variant. For a thorough evaluation, Table \ref{tab:distill} presents both the original LLaMA 3.1, trained on 15T tokens, and the {\textsc BLT} constructed from LLaMA 3. In our findings, the {\textsc BLT} model shows a minor decrease in accuracy on MMLU and HumanEval but demonstrates a more notable reduction on additional benchmarks. These results indicate that to maximize the benefits of the pre-trained model, additional research is necessary, especially in refining data composition and tuning hyperparameters.

\section{Ablations and Discussion}
\label{section:ablations}
Here, we present ablation studies that support the design decisions made for the {\textsc BLT} architecture, as well as choices regarding the patching scheme and the selection of hyper-parameters for the {\textsc BLT} model with 8B parameters, which was trained on the {\textsc BLT}-1T dataset.

\paragraph{Entropy Model Hyper-parameters} To evaluate how changes in entropy model size and the length of the context window impact scaling, we train several byte-level entropy transformer models, varying model capacities from 1 million to 100 million parameters and testing context windows spanning 64 to 512 tokens. The relationship between bits per byte (bpb) and training flops for these configurations is visualized in scaling law plots, based on our $400m$ and $1b$ {\textsc BLT} models trained with the Llama-2 dataset, as shown in \autoref{fig:entropy_ablation}. Our analysis indicates that larger entropy models and extended context windows both improve scaling results, although the benefit plateaus as the number of parameters exceeds 50 million.

\paragraph{Types of Patching}

We conduct ablation studies on the four patching methods described in Section~\ref{section:patching}: (1) Strided Patching with stride values of 4 and 6, (2) Whitespace-based patching, (3) BPE Tokenizer patching utilizing the Llama 3 tokenizer, and (4) Entropy-driven patching guided by a lightweight byte-level LLM.

Although dynamic patching shortens the sequences effectively, we standardize the sequence length across all patching methods to ensure a comparable context window. This approach ensures that, on average, each model is exposed to an equivalent number of bytes per sequence during both training and inference, thereby mitigating potential confounds associated with varying context sizes. The findings from these controlled experiments are presented in Figure~\ref{fig:scaling-compute-opt}. Notably, all alternative patching methods achieve superior performance relative to static patching, with space patching demonstrating performance nearly matching that of dynamic entropy-based patching.

In \autoref{tab:evals_ablation}, we report benchmark results for {\textsc BLT} models, contrasting tokenizer-based variants, models using space patching, and those employing entropy-based patching. All models were trained on the Llama 2 dataset for an optimized training duration~\citep{dubey2024llama}. While space patching is a more straightforward technique and avoids the need to compute entropies in real time during training, our analysis demonstrates that the benefits associated with entropy-based patching observed in scaling behaviors~(Section~\ref{section:scaling}) are also evident in downstream benchmark tasks.\footnote{The reported space patching outcomes derive from earlier experiments lacking cross-attention, though comparable patterns hold when cross-attention is included.}

\paragraph{Cross-Attention} 
In~\autoref{tab:crossattn_ablations_1b}, we examine the impact of introducing cross-attention modules at different positions within both the encoder and decoder components of {\textsc BLT}. For the encoder-side cross-attention, we consider three strategies for query initialization: 1) using an identical learned embedding across all global states, 2) employing a hash-based embedding derived from the bytes contained in each patch, and 3) applying a pooling operation to the encoder’s hidden states corresponding to the patch bytes at the specified encoder layer.

Our results indicate that incorporating cross-attention within the \textit{decoder} yields the greatest benefits. Within the encoder, marginal gains from cross-attention are observed, but these are limited to cases where queries are initialized with pooling. Furthermore, cross-attention proves especially advantageous on the Common-Crawl dataset and when employing larger patch sizes.

\paragraph{n-gram Hash Embeddings} 

We experiment with n-gram hash embedding vocabularies of sizes 0, 100K, 200K, and 400K, with corresponding results reported in Table~\ref{tab:ngram_ablations_1b}. Our analysis demonstrates that incorporating hash embeddings offers improvements across all tested domains, with particularly notable gains observed for Wikipedia and Github (showing a 0.04 bpb reduction versus 0.01 bpb after 15k steps at 8B scale). When increasing hashes from 300K to 500K at the 8B model size, the effect on performance is minor, leading to an approximate change of ~0.001 bpb at 15k steps. These results suggest hash embeddings are crucial for allowing {\textsc BLT} to achieve comparable performance with models relying on tokenization, but benefits plateau beyond 300K hashes. Moreover, our findings imply that hash embeddings and cross-attention offer largely additive benefits, enhancing performance on distinct datasets.

\paragraph{Local Model Hyperparamaters} Table \ref{tab:local_layers} presents an ablation study examining different configurations for the layer counts in the local encoder and decoder. The results indicate that, when utilizing hash n-gram embeddings, {\textsc BLT} achieves strong performance with a highly simplified encoder—specifically, one consisting of only a single layer—while employing a more complex decoder with additional layers.

\section{Related Work}
\label{section:related_work}

\textbf{Character-Level RNNs:} The task of Character Language Modeling has remained significant since the advent of early neural architectures~\citep{sutskever2011generating,mikolov2012subword,graves2013generating}, primarily because of its inherent ability to seamlessly handle out-of-vocabulary terms without the need for explicit back-off strategies. Building on this, \cite{kim2016character} introduced a model that exclusively processes character inputs by leveraging convolutional and highway networks, subsequently passing these representations to LSTM-based RNNs. Their approach was able to achieve parity with the then state-of-the-art RNN language models on English, and it surpassed them on languages with complex morphology—highlighting a key strength of character-level large language models. Similarly, \cite{kenter2018byte} applied byte-level LSTM architectures for machine comprehension, achieving superior performance relative to word-level counterparts, particularly for morphologically intricate languages such as Turkish and Russian. Continuing in this vein, \cite{zhang2015character} demonstrated that character-level convolutional models could deliver enhanced results over word-based systems for various classification challenges. \citet{chung2019hierarchical} explored hierarchical LSTM designs with boundary detectors at each layer, uncovering latent text hierarchies and further advancing character-level language modeling performance. Lastly, ByteNet as proposed by \cite{kalchbrenner2016neural} departs from attention mechanisms, opting instead for convolutional neural network layers applied directly to characters for the task of machine translation.

\textbf{Character-Level Transformers:} Transformer architectures leveraging attention mechanisms~\citep{vaswani2017attention}, alongside the adoption of subword tokenization techniques~\citep{sennrich-etal-2016-neural}, have yielded marked advancements in the effectiveness of neural models on language modeling and standard evaluation benchmarks. Nonetheless, employing word and sub-word representations introduces an inherent inductive bias regarding the abstractness at which these models operate. Seeking to merge the robust performance of transformer models with the encouraging outcomes in character-based language modeling, \citet{al2019character} explored the use of very deep transformer networks, further enhanced with auxiliary training objectives. Their approach enabled transformer-based architectures to exceed the performance of prior LSTM-driven character language models. Despite these gains, a substantial performance disparity in favor of word-level LLMs remained. Similarly, GPT-2~\citep{radfordgpt2} reported that when trained on large corpora such as the 1 Billion Word Benchmark, language models operating at the byte level still lagged behind those utilizing word-level representations.

Previous work by \cite{Choe2019BridgingTG} established that byte-level LLMs based on transformers can surpass subword-based models with similar parameter counts, yet these models demand considerably greater computational resources and extended training times. Likewise, \cite{el2020characterbert} introduced a variant of BERT (CharFormer), which constructs word representations through convolutions over character embeddings. While this model achieved superior results within the medical domain, it did so with increased computational costs. In parallel, \cite{clark2022canine} proposed CANINE, a 150M-parameter encoder-only model that processes raw character sequences. At its core, CANINE features a deep transformer stack, akin to the global model structure described in our work, coupled with a hybrid approach utilizing a local transformer and strided convolutions for downsampling characters. This architecture yields better performance on multilingual benchmarks compared to equivalent token-based encoder models such as mBERT. In the domain of byte-level encoder-decoder models, ByT5~\citep{xue2022byt5} explored implementations that omit patching mechanisms entirely. Although their approach demonstrated heightened robustness to input noise and was competitive with traditional tokenizer-based models using only a quarter of the data, the absence of patching required resource-intensive attention calculations across every byte, resulting in high computational overhead. Transitioning from subword to byte-level modeling notably inflates input sequence lengths, posing scalability obstacles for efficient training. Recently, leveraging the Mamba Architecture~\citep{gu2023mamba}, which preserves a constant-size memory state across extended contexts, \cite{wang2024mambabyte} developed a patchless, byte-level Mamba architecture. At the 350M-parameter scale, their model outperformed its byte-level transformer counterparts on several datasets in terms of bits-per-byte within a controlled FLOP budget.

\textbf{Patching-based approaches:} The strategic application of patching offers a way to reduce the large computational cost—measured in flops—characteristic of byte-level LLMs, without significantly sacrificing performance. Numerous studies have reported promising outcomes on relatively small-scale models and limited training byte budgets. For instance, \cite{nawrot-etal-2022-hierarchical} explored the use of static patching through downsampling and upsampling operations, introducing the hourglass transformer architecture, which demonstrates superior performance compared to other byte-level models at the 150M parameter scale. Building on this, \cite{nawrot-etal-2023-efficient} advance these results by implementing dynamic patching techniques, such as an end-to-end trainable boundary-predictor, a version of the predictor guided by specific tokenizers, and an entropy-driven patching mechanism akin to {\textsc BLT}. Their results indicate that this dynamic patching approach can surpass contemporaneous vanilla transformer models in language modeling tasks, particularly at a scale of 40M parameters trained on roughly 400M tokens. Meanwhile, \cite{lester2024training} probe the effectiveness of training on input sequences that have been compressed through arithmetic coding, achieving higher compression ratios than BPE. Employing the equal-info windows approach, they attain results superior to byte-level baselines on language modeling benchmarks, though performance still trails behind leading subword-based baselines.

This study is primarily influenced by, and most closely aligns with, MegaByte~\citep{yu2023megabyte}. MegaByte is a causal LLM employing only a decoder, where a predetermined, static approach is used to partition and merge representations to map bytes into patches. On the decoder end, a local model reconstructs bytes from these patches. Their results indicate that MegaByte is able to achieve comparable performance to models utilizing tokenizers, specifically at the 1B parameter scale when trained on approximately 400B bytes. In all our experimental evaluations, we perform ablation studies on MegaByte and determine that static patching does not outperform currently optimal tokenizer-based models that are trained with the same computational budget; moreover, we illustrate how {\textsc BLT} serves to close this performance gap. Similarly, \cite{slagle2024spacebyte} reach analogous conclusions about MegaByte, and propose enhancing the static patching method to include partitioning on whitespace and other space-like byte types, in addition to integrating a local encoder component. Their findings reveal that, in some areas (including Github and arXiv), these modifications yield better results than token-based transformer architectures, under compute-constrained conditions at the 1B parameter scale. In this paper, we also explore experiments with this particular model and observe that additional architectural advancements will be necessary for scaling byte-level models effectively and achieving parity with leading token-based architectures like Llama 3.

\section{Limitations and Future Work}

For this study, when making architectural decisions, we train our models for the number of steps deemed optimal for Llama~3~\citep{dubey2024llama}. Notably, the corresponding scaling laws were originally derived for BPE-based transformer models, and their application may result in less-than-ideal ratios of data to parameters for the {\textsc BLT} architecture. Determining more suitable scaling laws for {\textsc BLT}, which could reveal better scaling behavior, is left as a direction for future research. Moreover, as most of our experiments were performed with models containing up to 1B parameters, it remains an open question whether the most effective architectural choices might differ at scales of 8B parameters or higher, possibly yielding enhanced results at these larger model sizes.

Current transformer codebases prioritize optimization for architectures reliant on tokenizers. Although our experiments are theoretically matched in terms of floating point operations and utilize efficient tools like FlexAttention to process non-standard transformer layers, the wall-clock performance of our models might still lag behind tokenizer-based counterparts, indicating a potential for further improvements in runtime efficiency.

Although {\textsc BLT} employs an independently trained entropy model for the patching process, integrating the learning of the patching model into a unified end-to-end training regime presents a promising avenue for subsequent research. As discussed in Section \ref{sec:distillation}, we provide preliminary results that demonstrate encouraging progress in “byte-ifying” tokenizer-based models—including models like Llama 3 trained on datasets exceeding 10T tokens—by directly initializing and freezing the global transformer using their pretrained weights. Continued exploration along these lines may reveal approaches that not only preserve the advantages of bytefying but also enable performance improvements beyond those achievable by traditional tokenizer-based architectures, all without the necessity for training entirely new models from scratch.

\section{Conclusion}
\label{section:conclusion}

In this work, we introduce the Byte Latent Transformer (\textbf{BLT}), an innovative model architecture that moves away from the traditional reliance on predetermined vocabularies in large language models. {\textsc BLT} employs a novel, adaptive approach in which byte sequences are dynamically and learnably aggregated into patches, thereby allowing the distribution of computational effort to align with the underlying complexity of the input. This method results in marked gains in both computational efficiency and resilience to challenging input scenarios. Through an expansive scaling analysis, we establish that {\textsc BLT} achieves parity with tokenization-based architectures such as Llama 3 when scaled to model and dataset sizes of up to 8B parameters and 4T bytes, respectively, and can offer up to 50\% reductions in inference compute—albeit with only minor trade-offs in evaluation scores. Additionally, {\textsc BLT} introduces a novel axis for model scaling: it makes it possible to jointly grow both model size and patch size while maintaining a constant inference budget. This capability proves highly beneficial in the types of computing environments often seen in real-world applications. Notably, because {\textsc BLT} directly processes raw byte streams, it enhances the system’s capacity to manage rare or complex input cases, yielding superior robustness to noisy data and richer representations of sub-word elements. Collectively, these advances highlight {\textsc BLT} as a compelling successor to legacy tokenization schemes, delivering a highly scalable, resilient, and efficient infrastructure for future language modeling tasks.

\end{document}